\scanpage{52}
\mspage{- 35 -}
n\textsuperscript{o} 393

\deleted{21) Soient $A$ une algèbre de Banach unipollente, $\Delta$ une partie ouverte de $\handwritten{\mathbb C}$,
$\lambda\longrightarrow R(\lambda)$ une application de $\Delta$ dans $A$ telle que :
\[
R(\lambda)-R(\mu)=-(\lambda-\mu)R(\lambda)R(\mu)\tag*{(1)}
\]
quels que soient $\lambda$, $\mu\in\Delta$.}

\deleted{a) Soit $\lambda_{\unsure{o|0}}\in\Delta$. Pour tout $\lambda\in\Delta$ tel que $|\lambda-\lambda_o|\|R(\lambda)\|<1$,
on a :
\[
R(\lambda)=\sum_{n=0}^{\infty}(\lambda_o-\lambda)^n\,R(\lambda_o)^{n+1}.
\]}

\deleted{b) Pour tout $\lambda\in\Delta$, $\dfrac{d^n}{d\lambda^n}R(\lambda)=(-1)^n\,n!\,R(\lambda)^{n+1}$.}

\deleted{c) Si $R$ est holomorphe à l'infini, il existe $z$, $j$, $x\in A$ tels que $z^2=0$,
$j^2=j$, $zj=jz=0$, $x\in jAj$, et $R(\lambda)=z+jR(\lambda,x)$ pour $|\lambda|$ assez grand.
(Ecrire $R(\lambda)=\displaystyle\sum_{n=0}^{\infty}c_n\,\lambda^{-n}$ pour $|\lambda|>\lambda_o$, et exprimer que $R$ satisfait à (1).
On peut poser $c_o=z$, $c_1=j$, $c_2=x$).}

\deleted{d) Pour que $R$ soit la résolvante d'un élément de $A$, il faut et il suffit que
$R(\lambda_o)$ soit inversible pour un $\lambda_o\in\Delta$.}

\deleted{e) Pour tout $x\in A$, la fonction $\lambda\longrightarrow\displaystyle\sum_{n=0}^{\infty}(\lambda_o-\lambda)^n\,x^{n+1}$ satisfait à (1)
pour $|\lambda-\lambda_o|\ \|x\|<1$.}

\margin{\handwritten{\P}}
\deleted{22) a) Soient $B$ une algèbre de Banach unipollente, $w\in B$, $C>1$ et $\varepsilon>0$. Il
existe $\eta>0$ tel que les conditions $u\in B$, $\|u\|<C$, $\|uw-w\|<\eta$,
$0\leq\gamma\leq1/4C$ impliquent $\|yw-u\|<\varepsilon$, où $y=[(1-\gamma)e+\gamma u]^{-1}$.}

\deleted{b) Soit $A$ une algèbre de Banach. On suppose qu'il existe une constante $C$ telle
que, pour tout $\varepsilon>0$ et tous $x_1,\ldots,x_n\in A$, il existe $u\in A$ avec $\|u\|<C$,
$\|ux_1-x_1\|<\varepsilon,\ldots,\|ux_n-x_n\|<\varepsilon$. Pour tout $z\in A$, il existe $x$, $y\in A$
tels que $z=xy$, $y\in\overline{Az}$, $\|z-y\|<\delta$. (On peut supposer $C>1$. Soit $\gamma=1/4C$
Appliquant a) dans $\widetilde A$, déterminer par récurrence une suite $u_1$, $u_2$,\ldots\ dans $A$ telle}
% archive sheet number: 51

